% GENERATED FILE. Edit canonical lessons, then run npm run generate:books.
% canonical-source-hash: fnv1a64:dd1d6c896ddd02b4
% canonical-lessons: ML-C60-cow, ML-C60-goat, ML-C60-hen, ML-C60-crow, ML-C60-elephant

\chapter{Animals Around a House}
\label{ch:ml-animals-around-a-house}
\hlchaptermodality{60}

\begin{chapteropening}
\textbf{By the end of this chapter:} \emph{I can name a cow, a goat, a hen, a crow and an elephant in Malayalam.}

Complete ML-C60-elephant and say all five words in order.
\end{chapteropening}

\section[paśu]{\ml{പശു} (paśu) — a cow}

\label{lesson:ML-C60-cow}



\begin{quote}
Before the new run of words: what did \emph{puka} mean, and what did \emph{kanal} mean?
\end{quote}

\subsection*{You'll want to know: \ml{പശു}}
\textbf{\ml{പശു}} (\emph{paśu}) — \textquotedblleft{}a cow\textquotedblright{}.

Inherited from Sanskrit \emph{paśu}, which named any tethered beast a household kept, and which Malayalam has narrowed to the one animal a Kerala yard is arranged around.

The Sanskrit word has relatives a long way west. Latin \emph{pecu} meant a herd, and the English word \emph{fee} began as the same idea — the wealth you could drive to market on its own legs.

The milk-word \emph{pāl} comes out of a \ml{പശു}, and so does most of what a Kerala kitchen does next with it.

\begin{grammarlens}[title={What you've built}]
The first of five animals.
\end{grammarlens}

\subsection*{Your turn}
Say these aloud:

\begin{itemize}
\raggedright
  \item \emph{paśu}
  \item it once more, slowly
  \item \emph{paśu}, then \emph{pāl}, and say which one gives the other
\end{itemize}

\begin{tcolorbox}[breakable,colback=teal!4,colframe=teal!35!black,title={Before you move on}]
What does \ml{പശു} mean? (\textquotedblleft{}a cow\textquotedblright{}.) Say it once more.
\end{tcolorbox}
\section[āṭŭ]{\ml{ആട്} (āṭŭ) — a goat}

\label{lesson:ML-C60-goat}



\begin{quote}
Before the new one: what did \emph{paśu} mean?
\end{quote}

\subsection*{You'll want to know: \ml{ആട്}}
\textbf{\ml{ആട്}} (\emph{āṭŭ}) — \textquotedblleft{}a goat\textquotedblright{}.

Inherited, matching Tamil \emph{āṭu}, and it covers the animal at any age or either sex — a second word goes in front when the difference matters.

It sits one letter away from \emph{nāṭŭ}, the home-country word. Take the \emph{n} off the front of that one and you have this one, so the ear has to catch the very start of the word rather than the end of it.

A goat is what a household keeps when it has no room for the larger animal: it eats what it finds and it climbs what it likes.

\begin{grammarlens}[title={What you've built}]
Two.
\end{grammarlens}

\subsection*{Your turn}
Say these aloud:

\begin{itemize}
\raggedright
  \item \emph{āṭŭ}
  \item it once more, slowly
  \item \emph{nāṭŭ}, then \emph{āṭŭ}, and say which of the two is an animal
\end{itemize}

\begin{tcolorbox}[breakable,colback=teal!4,colframe=teal!35!black,title={Before you move on}]
What does \ml{ആട്} mean? (\textquotedblleft{}a goat\textquotedblright{}.) Say it once more.
\end{tcolorbox}
\section[kōḻi]{\ml{കോഴി} (kōḻi) — a hen}

\label{lesson:ML-C60-hen}



\begin{quote}
Before the new one: what did \emph{āṭŭ} mean?
\end{quote}

\subsection*{You'll want to know: \ml{കോഴി}}
\textbf{\ml{കോഴി}} (\emph{kōḻi}) — \textquotedblleft{}a hen\textquotedblright{}.

Native Dravidian, matching Tamil \emph{kōḻi}, and carrying the \ml{ഴ} that Malayalam and Tamil kept where their sisters let it go — the same letter sitting inside \emph{puḻa}, the river, and \emph{vaḻi}, the way.

One word does for hen and cock alike. Where a speaker needs to separate them, a second word goes in front.

It is also the clock of a Kerala house: a \ml{കോഴി} wakes the yard well before anyone has thought about the hour.

\begin{grammarlens}[title={What you've built}]
Three.
\end{grammarlens}

\subsection*{Your turn}
Say these aloud:

\begin{itemize}
\raggedright
  \item \emph{kōḻi}
  \item it once more, slowly
  \item \emph{kōḻi}, then \emph{puḻa}, and say which sound the two have in common
\end{itemize}

\begin{tcolorbox}[breakable,colback=teal!4,colframe=teal!35!black,title={Before you move on}]
What does \ml{കോഴി} mean? (\textquotedblleft{}a hen\textquotedblright{}.) Say it once more.
\end{tcolorbox}
\section[kākka]{\ml{കാക്ക} (kākka) — a crow}

\label{lesson:ML-C60-crow}



\begin{quote}
Before the new one: what did \emph{kōḻi} mean?
\end{quote}

\subsection*{You'll want to know: \ml{കാക്ക}}
\textbf{\ml{കാക്ക}} (\emph{kākka}) — \textquotedblleft{}a crow\textquotedblright{}.

An imitation word: the bird is named for the noise it makes, and the same trick produced the crow-word in each of the sister languages. Once you have heard one you have the etymology.

Of all the birds around a house this is the one a household deals with daily. Cooked rice — \emph{cōṟu} — set out on a step at the back is put there for the \ml{കാക്ക}, and it arrives within the minute.

Kerala also gives the bird a job at a funeral rite, where a ball of rice is set out for the ancestors and the crow that takes it is the sign the offering landed.

\begin{grammarlens}[title={What you've built}]
Four.
\end{grammarlens}

\subsection*{Your turn}
Say these aloud:

\begin{itemize}
\raggedright
  \item \emph{kākka}
  \item it once more, slowly
  \item \emph{kākka}, then \emph{kōḻi}, and say which of the two is kept and which turns up
\end{itemize}

\begin{tcolorbox}[breakable,colback=teal!4,colframe=teal!35!black,title={Before you move on}]
What does \ml{കാക്ക} mean? (\textquotedblleft{}a crow\textquotedblright{}.) Say it once more.
\end{tcolorbox}
\section[āna]{\ml{ആന} (āna) — an elephant}

\label{lesson:ML-C60-elephant}



\begin{quote}
Before the new one: what did \emph{kākka} mean?
\end{quote}

\subsection*{You'll want to know: \ml{ആന}}
\textbf{\ml{ആന}} (\emph{āna}) — \textquotedblleft{}an elephant\textquotedblright{}.

Inherited, and a worn-down cousin of Tamil \emph{yāṉai}: the same Dravidian word, with the front of it rubbed away over the centuries until two letters were left.

You met a piece of this animal before you met the animal. \emph{Kompŭ}, the branch-word, already covered the tusk of an \ml{ആന} — one picture for the branch, the horn and the tusk, and no separate word needed for the third.

In Kerala the \ml{ആന} belongs to a festival before it belongs to a forest. A caparisoned elephant standing in front of a temple is what most Malayalis picture on hearing the word.

\begin{grammarlens}[title={What you've built}]
Five, and the yard is full: the cow, the goat, the hen, the crow, and the one that walks in a festival.
\end{grammarlens}

\subsection*{Your turn}
Say these aloud:

\begin{itemize}
\raggedright
  \item \emph{āna}
  \item it once more, slowly
  \item all five in order — \emph{paśu}, \emph{āṭŭ}, \emph{kōḻi}, \emph{kākka}, \emph{āna}
\end{itemize}

\begin{tcolorbox}[breakable,colback=teal!4,colframe=teal!35!black,title={Before you move on}]
What does \ml{ആന} mean? (\textquotedblleft{}an elephant\textquotedblright{}.) Say it once more.
\end{tcolorbox}
